% !TeX root = Report.tex
\chapter*{Заключение}
\addcontentsline{toc}{chapter}{Заключение}

В работе рассматривалась задача сравнения алгоритмов предобработки изображения
по точности последующего оценивания параметров прямой методом Хафа. В качестве
основной модели использовалось изображение одной прямой на темном фоне с
аддитивным гауссовским шумом. Были рассмотрены как классические методы
шумоподавления, а именно медианный фильтр, фильтр Винера и \texttt{BM3D}, так и
алгоритм предобработки, основанный на применении \texttt{CSSA} после
дискретного преобразования Фурье.

В первой главе были описаны используемые алгоритмы и зафиксированы свойства
\texttt{CSSA}, существенные для дальнейшего анализа. Во второй главе
исследован механизм работы варианта \texttt{CSSA row.row} в случае одной прямой.
Показано, что в однокомпонентной постановке данный метод естественно
интерпретируется как выделение доминирующей компоненты сигнала в каждой строке
после преобразования Фурье. Вместе с тем установлено, что качество такой
реконструкции существенно зависит от уровня шума, расположения прямой на
дискретной сетке и правильности выбора компоненты. При высоком уровне шума
возможны эффекты растекания и смещения максимума, а компоненте сигнала может
соответствовать уже не первая, а одна из следующих компонент разложения.

В третьей главе проведено численное сравнение методов предобработки по точности
оценок параметров $\rho$ и $\theta$. Полученные результаты показывают, что
наиболее устойчивым в рассматриваемой постановке оказывается квантильный метод
\texttt{Quantile 0.9}. Его преимущество связано с тем, что после такой
предобработки множество активных пикселей оказывается существенно менее
плотным, чем после \texttt{Median}, \texttt{Wiener}, \texttt{BM3D} и
\texttt{CSSA row.row}. Для метода Хафа это обстоятельство принципиально,
поскольку точность оценивания определяется не столько визуальной гладкостью
изображения, сколько геометрией множества активных пикселей, участвующих в
голосовании. Показано также, что \texttt{CSSA row.row} может давать точные
результаты для части конфигураций, однако в некоторых конфигурациях
проявляется его неустойчивость, согласующаяся с выводами второй главы о
неправильном выборе компоненты и смещении максимума уже на уровне отдельных
строк. Дополнительный эксперимент показал, что уменьшение шага дискретизации
параметров преобразования Хафа приводит к более точному оцениванию параметров
прямой.

Таким образом, проведенное исследование показывает, что алгоритмы
предобработки следует сравнивать не только по качеству визуального
шумоподавления, но и по тому, насколько хорошо результат согласуется с
последующим методом оценивания. В рассматриваемой задаче именно такая связка
``предобработка --- преобразование Хафа'' оказывается определяющей.

Дальнейшее развитие работы может идти в нескольких направлениях. Во-первых,
представляет интерес автоматический выбор числа компонент в реконструкции
\texttt{CSSA}. Во-вторых, необходимо исследовать автоматический выбор
информативной компоненты в тех случаях, когда компоненте сигнала не
соответствует первое сингулярное число. В-третьих, естественным продолжением
является переход к случаю нескольких прямых на изображении, где задача выбора
числа компонент становится особенно существенной.
